\vspace{-2mm}
\section{Related Work}
\label{sec:related}
\vspace{-2mm}
We briefly discuss related work here and give an extended discussion of the related work in Appendix~\ref{appdx:extended_rw}.
%

Traditional RNNs are difficult to scale  due to the nonlinear dependencies between the  hidden states and expensive matmul-based sequential hidden state updates. Linear RNNs/State-Space Models (SSMs)/Transformers eliminate nonlinear dependencies, making training parallelizable along the temporal dimension \citep{parallel-martin, s4, s5}. Such models have been the focus of much recent work as a competitive sub-quadratic alternative to the Transformer architecture \citep{rwkv, Gu2023MambaLS,HGRN,qin2023scaling,sun2023retentive,pretrain_wo_attn}.




\begin{figure}[t] 
\vspace{-3mm}
		\centering 
		\resizebox{1\columnwidth}{!}{  
\begin{tikzpicture}
\scalefont{1.5} %
\begin{axis}[
    name=plot1,
    title={Training throughput},
    %
        xlabel={Training length/Batch size},
    ylabel={Tokens per second (Kt/s)},
    major x tick style=transparent,
    ybar=2*\pgflinewidth,
    bar width=8pt,
    x tick label style={rotate=0, anchor=center},
    symbolic x coords={2048/8,4096/4,8192/2,16284/1},
    xticklabel style={yshift=-2mm}, 
    ymin=0,
    xtick=data,
    tick label style={font=\large},
    label style={font=\large},
    %
    legend style={at={(7,-1.6)},   
                anchor=north,legend columns=3,
                column sep=0.3cm,
                font=\normalsize,
                row sep=-0.15cm
                },     
    legend cell align=left,
    %
    ymajorgrids=true,
    grid style=dashed,
    %
    %
    enlarge x limits=0.25,
    %
    point meta=explicit symbolic,
    legend to name={mylegend},
]
%
\addplot[draw=bred,thick,fill=bred,fill opacity=0.6] coordinates { 
  (2048/8, 51.3)[\textcolor{bred!60}]
  (4096/4, 46.7)[\textcolor{bred!60}]
  (8192/2, 38.7)[\textcolor{bred!60}]
   (16284/1, 29.1)[\textcolor{bred!60}]
};

%
\addplot[draw=bblue,thick,fill=bblue,fill opacity=0.6] coordinates {
  (2048/8, 22.8)[\textcolor{bblue!60}{1408}]
  (4096/4, 22.8)[\textcolor{bblue!60}{352}]
  (8192/2, 22.8)[\textcolor{bblue!60}{1056}]
  (16284/1, 26.0)[\textcolor{bblue!60}{1056}]
};

\addplot[draw=bcyan,thick,fill=bcyan,fill opacity=0.6] coordinates {
  (2048/8, 43.8)[\textcolor{bcyan!60}{1056}]
  (4096/4, 43.5)[\textcolor{bcyan!60}{352}]
  (8192/2, 43.2)[\textcolor{bcyan!60}{1056}]
  (16284/1, 41.1)[\textcolor{bcyan!60}{1056}]
  };


%
%
%
%
%
%

\legend{Transformer++, Mamba, GLA}

\end{axis}
\begin{axis}[
 at={(plot1.south east) },
    title={GPU memory usage},
    xlabel={Training length/Batch size},
    ylabel near ticks, yticklabel pos=right,
    ylabel={Gigabyte (GB)},
    major x tick style=transparent,
    ybar=2*\pgflinewidth,
    bar width=8pt,
    x tick label style={rotate=0, anchor=center},
    symbolic x coords={2048/8, 4096/4, 8192/2, 16284/1},
    xticklabel style={yshift=-2mm},
    %
    %
    xtick=data,
    ymin=0,
      tick label style={font=\large},
  label style={font=\large},
    %
    %
    ymajorgrids=true,
    grid style=dashed,
    %
    %
    enlarge x limits=0.25,
    %
    point meta=explicit symbolic
]
\addplot[draw=bred,thick,fill=bred,fill opacity=0.6] coordinates { 
  (2048/8,33)[\textcolor{bred!60}]
  (4096/4, 33)[\textcolor{bred!60}]
  (8192/2, 33)[\textcolor{bred!60}]
   (16284/1, 33)[\textcolor{bred!60}]
};

%
\addplot[draw=bblue,thick,fill=bblue,fill opacity=0.6] coordinates {
  (2048/8, 36.0)[\textcolor{bblue!60}{1408}]
  (4096/4, 36.0)[\textcolor{bblue!60}{352}]
  (8192/2, 36.0)[\textcolor{bblue!60}{1056}]
  (16284/1, 36.0)[\textcolor{bblue!60}{1056}]
};

\addplot[draw=bcyan,thick,fill=bcyan,fill opacity=0.6] coordinates {
  (2048/8, 37.0)[\textcolor{bcyan!60}{1056}]
  (4096/4, 37.0)[\textcolor{bcyan!60}{352}]
  (8192/2, 37.0)[\textcolor{bcyan!60}{1056}]
  (16284/1, 37.0)[\textcolor{bcyan!60}{1056}]
  };
\end{axis}
\ref{mylegend}

\end{tikzpicture}
		}
     \vspace{-6mm}
		%
		\caption{Training throughput and  memory footprint on an H100.} 
\label{fig:tps}
	\end{figure}




Data-dependent decay rates have always been regarded important for  RNNs  \citep{DBLP:journals/neco/GersSC00,unreasonable-forget-gate}. Typical forget gate values depend on both the previous hidden state and the current input. However \citet{parallel-martin} suggest that forget gate values should depend solely on the current inputs to enable parallel training. This simple strategy has been shown to be effective in moderate-scale experiments conducted by HGRN \citep{qin2023scaling}.  RWKV-v6 \citep{peng2024eagle} and Mamba \citep{Gu2023MambaLS} also use data-dependent decay rates that are reminiscent of forget gates. In the context of linear Transformers, \citet{peng2021random} employ a coarse-grained position-wise forget gate, while \citet{mao-2022-fine} and \citet{gatedloop} use a more fine-grained forget gate. 
%

RNNs rely on fixed-dimensional hidden states to encode their entire history. The hidden state dimension  serves as a proxy for memory capacity and thus significantly influences their expressive power. 
Linear Transformers expand the hidden dimension of RNNs via the outer-product parameterization, as discussed \S\ref{subsec:background-lin}. Linear SSMs on the other hand expand their hidden dimension via a single-input-single-output (SISO) strategy. 
Without data-dependent SSM parameters, this can be done efficiently during training via the Fast Fourier Transform (FFT). However, with data-dependent SSM parameters, FFT-based training is not possible, and thus \citet{Gu2023MambaLS} implements a custom CUDA kernel to train a selective state-space model using the parallel scan algorithm \citep{s5}. To fit all the hidden states into SRAM, they can only afford an expansion rate up to 16. In contrast our hardware-aware training algorithm provides an alternative, efficient approach for expanding the hidden dimension to a wider range, which we have shown useful in recall-intensive tasks.


\vspace{-2mm}
%
%





